\section{Generalization to Non-Code Settings}
\label{app:non-code}
\subsection{Few-shot antonym example}
The saved exploratory run in \texttt{19\_goalscope\_out\_of\_domain.ipynb} uses six demonstrations: enter/exit, feasible/infeasible, manned/unmanned, fuzzy/smooth, arise/settle, and far/near. The prompt ends at \texttt{Q: mercy} followed by \texttt{A:}. Its continuation starts with ``cruelty'' and then generates further question--answer pairs. The code-localized 122-head set supplies norm-relative additive readouts with $\alpha=.1$, injection layers 35--45, and repeated steering during readout generation.

The probe uses the system instruction ``You are a helpful assistant.'' and the user instruction ``You are answering a series of questions that all follow the same hidden rule. In one short phrase, what task are you performing?'' This is a domain-adapted probe, not the original code-oriented prompt. Saved token indices refer to the notebook's tokenized prompt plus continuation.

\begin{table}[h]\centering\small
\begin{tabular}{rll}\toprule
Position & Source token & Readout\\\midrule
57 & newline after cruelty & Identifying antonyms.\\
60 & light & Identifying patterns.\\
70 & A before slow & Identifying antonyms.\\
71 & colon before slow & Decoding or interpreting slower actions.\\
\bottomrule\end{tabular}
\caption{Selected qualitative readouts from one saved antonym run, including a vague description and a mismatch. These are not a quantitative evaluation sample.}
\end{table}

\subsection{Other exploratory domains}
Recipe, travel, and narrative prompts use adapted verbalization questions. These settings require separate decisions about the granularity and correctness of the next step; they are not included in the code results or claimed as a validated general-purpose planning readout. A systematic evaluation remains to be completed.
